% ===================== AgentErrorBench =====================
\subsection{AgentErrorBench (AgentDebug)}\label{ds:agenterrorbench}
\naturalbadge\ \okbadge\ \lbl{step + module + error type}\ \lbl{single-agent}\ \lbl{ALFWorld / GAIA / WebShop}

\begin{identity}
Paper & \emph{Where LLM Agents Fail and How They Can Learn From Failures} (ICML listing: \emph{AgentDebug: \dots}) --- Zhu, Liu, Li, Tian, Yang, Zhang, Han, Xie, Cui, Zhang, Ma, Yu, Ramesh, Wu, Liu, Lu, Zou, You (UIUC U~Lab, with Stanford / AMD), ICML 2026, arXiv:2509.25370 \cite{agentdebug2026}. (The id 2509.23071 sometimes quoted for this paper is a different paper, EviPath.)\\
Links & \href{https://github.com/ulab-uiuc/AgentDebug}{GitHub \code{ulab-uiuc/AgentDebug}} (code, no data) \textperiodcentered\ \href{https://drive.google.com/drive/folders/1bQe6dQA85pktT63YnKIKJDTVaH3O3Vpu}{Google Drive folder} (official data, linked from the README) \textperiodcentered\ \href{https://huggingface.co/datasets/davide221/agenterrorbench}{HF \code{davide221/agenterrorbench}} (third-party parquet mirror) \textperiodcentered\ re-packaged as \code{agentdebugbench} in \href{https://github.com/THU-KEG/TrajDebug}{\code{THU-KEG/TrajDebug}}\\
License & code MIT (``U Lab @UIUC''); no license file on the Drive data (ALFWorld / WebShop / GAIA terms apply); paper CC BY-NC-SA 4.0\\
Status & downloaded in full: the Drive folder yielded 196 of 200 trajectories (98 ALFWorld / 48 GAIA / 50 WebShop) + 3 label files, the 4 missing ones were filled from the HF mirror; 124 MB; \code{data/agenterrorbench/\{gdrive, hf\_davide221, repo\_AgentDebug\}/}. The camera-ready additionally announces \textbf{AgentDebug-Traj2k} ($\approx$2\,k proposer--reviewer-annotated trajectories) and the \textbf{AgentDebugger-4B} detector: \nobadge\ --- no artefact on GitHub, HF or Drive as of 2026-10-16.\\
\end{identity}

\begin{multicols}{2}
\paragraph{Description} AgentErrorBench is the test-bed of the \emph{AgentDebug} framework: given a failed single-agent rollout, predict (i) the step of the first critical (root-cause) error, (ii) the agent \emph{module} that produced it and (iii) the fine-grained error type from the accompanying \emph{AgentErrorTaxonomy}; AgentDebug then re-rolls the agent from that step with corrective feedback. The authors ran their own modular ReAct-style agent (each turn emits \code{<memory>}, \code{<reflection>}, \code{<plan>}, \code{<action>} blocks; 30-step cap) with GPT-4o, Llama-3.3-70B-Turbo and Qwen3-8B on three environments, collected $>$500 natural failures, derived the taxonomy from them and curated \textbf{200 representative failed trajectories}:
\begin{itemize}
\item \textbf{ALFWorld} (100): embodied household text game; every released run hits the 30-step cap.
\item \textbf{WebShop} (50): e-commerce browsing with \code{search[]} / \code{click[]} actions.
\item \textbf{GAIA} (50): tool-using research assistant; the shortest runs (1--30 steps).
\end{itemize}
\textbf{Annotation}: ten graduate annotators labelled decision steps with taxonomy types and marked the minimal set of root-cause steps; guidelines were refined over three pilot rounds, a shared subset was double-annotated and adjudicated; Cohen's $\kappa$ = 0.55 across modules (paper). In the release every trajectory carries exactly \emph{one} annotated step (= \code{critical\_failure\_step}) with one module, one type and a free-text reasoning; no repair text is included. \textbf{Taxonomy} (5 modules / 17 types): Memory (over-simplification, hallucination, retrieval failure), Reflection (progress misjudgement, outcome misinterpretation, causal misattribution, hallucination), Planning (constraint ignorance, impossible action, inefficient planning), Action (planning--action disconnect, format error, parameter error) and System (step-limit exhaustion, tool execution error, LLM limit, environment error).

\paragraph{Metrics and baselines} \textbf{Step} accuracy (exact step), \textbf{Step+Module} and \textbf{All-Correct} (step + module + type), detector backbone GPT-4.1 at temperature 0; cells are Step / Step+Module / All-Correct in \% (paper Table 1):

\begin{center}\footnotesize\setlength{\tabcolsep}{2.5pt}
\begin{tabular}{@{}lcccc@{}}\toprule
Method & ALFWorld & WebShop & GAIA & Avg\\\midrule
Direct prompting & 28 / 14 / 1 & 30 / 6 / 0 & 26 / 10 / 0 & 28.0 / 10.0 / 0.3\\
Brute force & 10 / 5 / 0 & 8 / 0 / 0 & 18 / 8 / 0 & 12.0 / 4.3 / 0.0\\
Binary search & 20 / 6 / 1 & 14 / 8 / 0 & 22 / 10 / 0 & 18.7 / 8.0 / 0.3\\
AgentDebug & \textbf{35 / 28 / 21} & \textbf{42 / 22 / 14} & \textbf{58 / 44 / 38} & \textbf{45.0 / 31.3 / 24.3}\\\bottomrule
\end{tabular}
\end{center}
Brute-force re-rolling (substitute a corrected action at each step until the run succeeds) is \emph{worse} than prompting. TrajDebug later re-evaluates the same 200 runs: direct prompting 9--19\% (ALFWorld), 20--35\% (GAIA), 14--32\% (WebShop) step accuracy; TrajDebug 26.0 / 38.8 / 26.0. Downstream, AgentDebug feedback plus up to five re-rolls yields up to 26\% relative task-success gain (paper).

\paragraph{Schema} \code{Label/\{alfworld,gaia,webshop\}\_labels.json}: list of \{\code{trajectory\_id}, \code{LLM}, \code{task\_type}, \code{critical\_failure\_step} (1-based agent turn $k$ = message index $2k{-}1$), \code{critical\_failure\_module}, \code{step\_annotations} = [\{\code{step}, \code{<module>}: \{\code{failure\_type}, \code{reasoning}\}\}]\}. \code{Original\_Failure\_Trajectory/<Env>/<id>.json}: \code{messages} (alternating \code{user} = task prompt + observation + admissible actions, or tool/web output on GAIA, and \code{assistant} = the tagged blocks, with \code{<think>} for Qwen3) and \code{metadata} (\code{model}, \code{steps}, \code{won} = False, \code{env\_id}, \dots). The HF mirror flattens this to one parquet row per trajectory (\code{full\_trajectory} as a JSON string, \code{failure\_types[]}, \code{num\_steps}) with an uploader-made 139 / 31 / 30 split; the official release has no split.
\end{multicols}

\paragraph{Statistics} computed from the downloaded data.\par\nobreak\vspace{1pt}
\noindent\begin{minipage}[t]{0.56\linewidth}\vspace{0pt}\small\setlength{\tabcolsep}{4pt}
\begin{tabular}{@{}lrrrr@{}}\toprule
 & All & ALFWorld & GAIA & WebShop\\\midrule
trajectories (all failed) & 200 & 100 & 50 & 50\\
agent steps: mean / med. & 25.1 / 30 & 30 / 30 & 14.3 / 10 & 25.9 / 30\\
agent steps: min -- max & 1 -- 30 & 30 -- 30 & 1 -- 30 & 4 -- 30\\
annotated steps / trajectory & 1 & 1 & 1 & 1\\
critical step: mean / med.\ (max) & 7.9 / 5 (30) & --- & --- & ---\\
rel.\ error position: mean / med. & 0.38 / 0.23 & --- & --- & ---\\
chars / trajectory: mean / med. & 378\,k / 198\,k & --- & --- & ---\\\bottomrule
\end{tabular}

\smallskip Backbones: GPT-4o 81 (50 / 16 / 15 per environment), Llama-3.3-70B-Turbo 57 (23 / 17 / 17), Qwen3-8B 62 (27 / 17 / 18). 155 of 200 runs are exactly 30 steps long (the cap), 32 have $\le$9. The critical step lies in the first fifth of the run in 78 cases and in the last fifth in 38; the paper's claim that failures ``cluster in mid-trajectory steps 6--15'' does not match the released labels (median 5). Blamed module: planning 39\%, reflection 20\%, memory 19\%, system 11\%, action 11\% --- the paper's Figure~3 shows memory and reflection on top, presumably counting all per-step tags from the annotation study rather than the single released label. 27 labels have an empty \code{failure\_type}. Trajectories are huge (max 5.8\,M characters) because WebShop/GAIA observations repeat the whole prompt and page each step.
\end{minipage}\hfill
\begin{minipage}[t]{0.42\linewidth}\vspace{0pt}
\includegraphics[width=\linewidth]{../figures/cat_agenterrorbench.pdf}
\end{minipage}

\paragraph{Example trajectory} \code{Qwen3-8B\_014\_chat\_b000\_t00\_e14-1d03d2d8} (WebShop, Qwen3-8B; 4 agent steps, 26\,k characters; each step = observation + the agent's tagged response). The agent searches, opens a shirt in the right colour and size, and clicks \code{buy now} without ever checking the ``machine wash'' requirement; the label is \code{reflection.progress\_misjudge} at step 4, the final step.
% auto-generated by scripts/render_traj.py -- do not edit
\trajtask{Qwen3-8B\_014\_chat\_b000\_t00\_e14-1d03d2d8}{You are an expert agent operating in the WebShop e-commerce environment.\newline Your task is:\allowbreak{} Find me machine wash men's t-shirts with long sleeve with color:\allowbreak{} mossy oak bottomland,\allowbreak{} and size:\allowbreak{} 5x-large big,\allowbreak{} and price lower than 50.00 dollars\newline Your current observation is:\allowbreak{} 'Search'\newline Your admissible actions of the current situation are:\allowbreak{} 'search[<your query>]',\allowbreak{}\newline 'click[search]',\allowbreak{}.\par Please begin by analyzing the situation and planning your approach:\allowbreak{}\par <plan>\newline Plan the next step:\allowbreak{}\newline - Given what I've learned,\allowbreak{} what should\newline [...] (71 chars omitted) [...]\newline n?\allowbreak{}\newline - What do I expect this action to achieve?\allowbreak{}\newline </\allowbreak{}plan>\par <action>\newline Finally,\allowbreak{} choose ONE admissible action for the current step and choose it within 'search[<your query>]',\allowbreak{}\newline 'click[search]',\allowbreak{}.\newline </\allowbreak{}action>}
\begin{trajstep}{1}{assistant (single ReAct-style agent)}{agA}
Observation (environment/\allowbreak{}user message 0):\allowbreak{}\newline (task prompt + initial observation;\allowbreak{} see `task`)\par Agent output (assistant message 1):\allowbreak{}\newline <think>\newline Okay,\allowbreak{} let's tackle this query. The user wants machine wash men's t-shirts with long sleeves,\allowbreak{} color mossy oak bottomland,\allowbreak{} size 5x-large,\allowbreak{} and under \$50. First,\allowbreak{} I need to parse all the details.\par The main challenge is the color:\allowbreak{} "mossy oak bottomland." That's a specific shade,\allowbreak{} maybe a military or outdoor-themed color. I should check if that's a standard color name or if it's a brand-specific term. If it's not common,\allowbreak{} the search might not find exact matches. But since the user specified it,\allowbreak{} I have to include it.\newline [...] (1,\allowbreak{}801 chars omitted) [...]\newline k bottomland color,\allowbreak{} 5x-large size,\allowbreak{} and price under \$50. This action will efficiently narrow down results to meet all requirements.\newline </\allowbreak{}plan>\par <action>\newline search[men's long sleeve t-shirt mossy oak bottomland 5x large under 50 dollars machine wash]\newline </\allowbreak{}action>
\end{trajstep}
\begin{trajstep}{2}{assistant (single ReAct-style agent)}{agA}
Observation (environment/\allowbreak{}user message 2):\allowbreak{}\par You are an expert agent operating in the WebShop e-commerce environment.\newline Your task is to:\allowbreak{} Find me machine wash men's t-shirts with long sleeve with color:\allowbreak{} mossy oak bottomland,\allowbreak{} and size:\allowbreak{} 5x-large big,\allowbreak{} and price lower than 50.00 dollars\newline Prior to this step,\allowbreak{} you have already taken 1 step(s). Below is a compact summary of all steps:\allowbreak{} - SearchQuery:\allowbreak{} men's long sleeve t-shirt mossy oak bottomland 5x large under 50 dollars machine wash\newline - PagesVisited:\allowbreak{} unknown\newline - RelevantProducts:\allowbreak{} unknown\newline - Selections:\allowbreak{} none\newline - IrrelevantSummary:\allowbreak{} none\newline You are now at step 2 and your current observation is:\allowbreak{} 'Back to Search' [SEP\newline [...] (2,\allowbreak{}605 chars omitted) [...]\newline ,\allowbreak{} so it's a logical candidate. If it doesn't match the color,\allowbreak{} the user can proceed to the next product or page. This action helps narrow down the options and ensures the user doesn't miss potential matches.\newline </\allowbreak{}plan>\par <action>\newline click[b00o30jldk]\newline </\allowbreak{}action>
\end{trajstep}
\begin{trajstep}{3}{assistant (single ReAct-style agent)}{agA}
Observation (environment/\allowbreak{}user message 4):\allowbreak{}\par You are an expert agent operating in the WebShop e-commerce environment.\newline Your task is to:\allowbreak{} Find me machine wash men's t-shirts with long sleeve with color:\allowbreak{} mossy oak bottomland,\allowbreak{} and size:\allowbreak{} 5x-large big,\allowbreak{} and price lower than 50.00 dollars\newline Prior to this step,\allowbreak{} you have already taken 2 step(s). Below is a compact summary of all steps:\allowbreak{} ```\newline SearchQuery:\allowbreak{} men's long sleeve t-shirt mossy oak bottomland 5x large under 50 dollars machine wash\newline PagesVisited:\allowbreak{} Page 1\newline RelevantProducts:\allowbreak{}\newline [B00O30JLDK] --- Legendary Whitetails Men's Non-Typical Long Sleeve T-Shirt --- \$10.52 to \$40.5 --- [Attrs:\allowbreak{} color=\allowbreak{}unknown,\allowbreak{} size=\allowbreak{}unknown,\allowbreak{}\newline [...] (2,\allowbreak{}605 chars omitted) [...]\newline filter and then check if '5x-large big' is available. This will narrow down the results to match the exact specifications,\allowbreak{} ensuring the product meets all criteria including price,\allowbreak{} color,\allowbreak{} and size.\newline </\allowbreak{}plan>\par <action>\newline click[mossy oak bottomland]\newline </\allowbreak{}action>
\end{trajstep}
\begin{faultystep}{4}{assistant (single ReAct-style agent)}{agA}
Observation (environment/\allowbreak{}user message 6):\allowbreak{}\par You are an expert agent operating in the WebShop e-commerce environment.\newline Your task is to:\allowbreak{} Find me machine wash men's t-shirts with long sleeve with color:\allowbreak{} mossy oak bottomland,\allowbreak{} and size:\allowbreak{} 5x-large big,\allowbreak{} and price lower than 50.00 dollars\newline Prior to this step,\allowbreak{} you have already taken 3 step(s). Below is a compact summary of all steps:\allowbreak{} SearchQuery:\allowbreak{} men's long sleeve t-shirt mossy oak bottomland 5x large under 50 dollars machine wash\newline PagesVisited:\allowbreak{} Page 1\newline RelevantProducts:\allowbreak{}\newline [B00O30JLDK] --- Legendary Whitetails Men's Non-Typical Long Sleeve T-Shirt --- \$10.52 to \$40.5 --- [Attrs:\allowbreak{} color=\allowbreak{}mossy oak bottomland,\allowbreak{} size\newline [...] (2,\allowbreak{}605 chars omitted) [...]\newline <plan>\newline The next step is to finalize the purchase by clicking 'buy now' to complete the transaction. This action directly achieves the task goal of acquiring the specified product at the correct price and size.\newline </\allowbreak{}plan>\par <action>\newline click[buy now]\newline </\allowbreak{}action>
\end{faultystep}
\trajannotation{assistant (single agent)}{4}{reflection.progress\_misjudge}{The agent didn't check if machine wash is met. The reflection is not able to recognize this correctly. It should check the product details to see if machine wash is met.}{—}{<think>\newline Okay,\allowbreak{} let's tackle this step by step. The user is looking for a machine wash men's long sleeve t-shirt in mossy oak bottomland color,\allowbreak{} size 5x-large,\allowbreak{} under \$50. From the history,\allowbreak{} t\newline [...] (1,\allowbreak{}855 chars omitted) [...]\newline at the correct price and size.\newline </\allowbreak{}plan>\par <action>\newline click[buy now]\newline </\allowbreak{}action>}{\textbf{critical\_failure\_module}: reflection; \textbf{step\_annotations}: [\{"step":\allowbreak{} 4,\allowbreak{} "reflection":\allowbreak{} \{"failure\_type":\allowbreak{} "progress\_misjudge",\allowbreak{} "reasoning":\allowbreak{} "The agent didn't check if machine wash is met. The reflection is not able to recognize this correctly. It should check the product details to see if machine wash is met."\}\}]; \textbf{step\_numbering}: 1-based index of assistant turns (message index 2k-1); \textbf{llm}: Qwen3-8B; \textbf{env}: webshop; \textbf{won}: False}


\paragraph{Notes} (i) Single-agent: ``who'' is always the one ReAct agent, ``when'' is a 1-based agent turn; the label adds a module and a 17-type error class, which is finer than Who\&When but coarser than AgentHallu's cause text. (ii) ALFWorld runs are uniformly capped at 30 steps, so step-limit exhaustion masks earlier root causes, and the label is often very early (median relative position 0.23). (iii) Label-quality quirks: module/type spellings vary (\code{plan} vs \code{planning}, \code{plan\_inefficient} vs \code{inefficient\_plan}, \code{Parameter\_error}, trailing spaces), 27/200 types are empty, one GAIA record (\code{GPT-4o\_003\_memory\_\dots}) labels step 3 of a 1-step trajectory, and only one step per trajectory is annotated although the paper describes multi-step annotation; $\kappa$ = 0.55 is the lowest agreement among the human-labelled sets here. (iv) Context budget: median 198\,k characters per trajectory, so detectors need truncation or compression (TrajDebug's motivation). (v) Traj2k and AgentDebugger-4B are announced in the ICML camera-ready only --- cite as ``not released''. (vi) Overlaps: the same 200 trajectories ship inside TrajDebug/TrajErrBench (\S\ref{ds:trajerrbench}) in a unified 0-based message schema, and that paper's taxonomy re-organises these five modules; Who\&When is the closest multi-agent sibling (GAIA), while ALFWorld and WebShop appear nowhere else in this survey.
